模式帮助我们应对复杂性。在单个软件组件的层面，你可以应用著名的经典 \emph{四人组（Gang of Four）}软件模式——即《设计模式：可复用面向对象软件的基础》（1994）一书中所描述的那些模式。然而，当我们站得更高，开始审视不同组件之间的架构，并研究组件如何交互与协同运行时，懂得何时以及如何应用架构与系统设计模式就变得至关重要，而这正是本章的重点。

尽管本章并不直接聚焦于 C++ 代码，但它讨论的是每一位有经验的 C++ 开发者最终都会面对、尤其是在处理复杂的分布式系统时会遇到的架构与系统设计问题。这里讨论的模式与技术无关，但对于追求可扩展性、容错性和生产就绪的 C++ 项目却高度相关。

适用于不同场景的这类模式数不胜数。事实上，想要全部了解它们，你需要读的可不止一本书。话虽如此，我们还是为本书挑选了几个模式，它们适合用来实现各种架构目标。

在本章中，我们将向你介绍几个与架构设计相关的概念和谬误。我们还会展示何时使用前述模式，以及如何设计易于部署的高质量组件。

本章将涵盖以下主题：

\begin{itemize}
\item 理解分布式系统的特性
\item 让你的系统容错且可用
\item 集成你的系统
\item 在大规模下实现性能
\item 设计数据存储
\item 部署你的系统
\item 管理你的 API
\end{itemize}

到本章结束时，你将懂得如何设计架构来提供几项重要的品质，例如容错性、可扩展性和可部署性。

\section{技术要求}

本章的代码需要以下工具来构建和运行：Docker 和 Docker Compose。

本章的源代码片段可在 \href{https://github.com/PacktPublishing/Software-Architecture-with-Cpp-2E/tree/main/Chapter04}{https://github.com/PacktPublishing/Software-Architecture-with-Cpp-2E/tree/main/Chapter04} 找到。

\section{理解分布式系统的特性}

\label{idx:261}
软件系统有很多种类型，每一种都适用于不同的场景、为不同的需求而构建，并基于不同的假设。一般来说，编写和部署一个传统的独立桌面应用，与编写和部署一个需要通过网络与其他众多服务通信的微服务，完全是两回事。

在本节中，我们将逐一介绍可用于部署软件的各种模型、人们在创建分布式系统时应避免的常见错误，以及人们为成功创建这类系统所做的某些权衡。

\subsection{不同的服务模型及其适用场景}

\label{idx:262}
\label{idx:263}
我们先从服务模型说起。在设计复杂的分布式系统时，服务模型界定了你需要自行管理多少基础设施，以及你能在多少现有构件之上进行构建。有时，你可能想利用现有的软件，而无须手动部署应用或备份数据，例如通过 API 使用 Google Drive 作为应用存储。另一些时候，你可以依赖现有的云平台（如 Google 的 App Engine）来部署你的解决方案，而无须操心语言运行时、负责管理程序执行的软件栈或数据库。如果你能自行决定以任意方式部署一切，那么你既可以借助云提供商的基础设施，也可以使用你自己公司的基础设施。下面的图示展示了不同的服务模型：

\myGraphic{0.920}{images/B18456_04_01.png}{图 4.1：服务模型}

下面我们来讨论这些不同的模型，以及每一种在哪些情况下有用。

\subsubsection{本地部署模式}

\label{idx:264}
\label{idx:265}
最经典的做法，也是前云计算时代唯一可行的做法，就是把一切都部署在你自己的本地环境中。你需要购买所需的全部硬件和软件，并确保它能提供满足你需求的足够容量。如果你在为一家初创公司工作，这可能会是一笔很大的前期成本。随着用户群的增长，你需要购买并配置更多资源，这样你的服务才能应对哪怕偶尔出现的负载高峰。这一切都意味着你必须预测解决方案的增长并主动采取行动，因为根本无法根据当前负载自动伸缩。

即便在云计算时代，本地部署依然有用，而且在现实中经常能见到。有时你处理的数据出于数据隐私或合规方面的原因，不应该、甚至不能离开公司的本地环境。另一些时候，你需要尽可能低的网络延迟，因而需要自己在地理上就近的数据中心来托管应用。有时你可能会核算成本，然后认定在你这种情况里，本地部署比云方案更便宜。最后但同样重要的是，你的公司可能已经有一座现成的数据中心可供使用。

\label{idx:266}
本地部署并不意味着你必须有一个单体系统。公司常常在本地部署自己的私有云。这有助于更好地利用现有基础设施来削减成本。你也可以把私有云方案与其他服务模型之一相结合，当你时不时需要额外容量时会很有用。这被称为\textbf{混合部署}，所有主流云提供商都提供这种方案，OpenStack 的 Omni 项目也提供了它。

混合云是私有云和公有云之间的折中。私有云通常能对数据安全与隐私进行细粒度控制，但成本更高、可扩展性也更有限；而公有云更具可扩展性、更具成本效益，但往往缺乏同等程度的控制力。

\subsubsection{基础设施即服务模式}

\label{idx:267}
\label{idx:268}
\label{idx:269}
说及其他模型，最基础的云服务模型叫做\textbf{基础设施即服务 }（\textbf{IaaS}）。它也是最接近本地部署的一种：你可以把 IaaS 理解为拥有一座虚拟数据中心的方式。顾名思义，云提供商把他们托管的基础设施的一部分提供给你，其中包含三类资源：

\begin{itemize}
\item 计算资源，例如虚拟机、容器或裸金属机器（不含操作系统）
\item 网络，除网络本身外，还包括 DNS 服务器、路由和防火墙
\item 存储，包括备份和恢复能力
\end{itemize}

提供所有软件仍然由你负责：操作系统、中间件以及你的应用。

\label{idx:270}
IaaS 可用于各种场景，例如托管网站（可能比传统虚拟主机更便宜）、提供存储（例如 Amazon 的 S3 和 Glacier 服务），以及支持高性能计算和大数据分析（需要巨大的计算能力）。有些公司用它来在需要时快速搭建和清除测试与开发环境。

使用 IaaS 而非本地基础设施，是一种测试新想法的经济高效的方式，同时还能省去配置所需的时间。

如果你的服务出现使用量高峰，例如在周末，你可能想利用云的自动伸缩能力：需要时扩容，之后再缩容以节省开支。

所有流行的云服务提供商都提供 IaaS 解决方案。

\begin{myNotic}{注意}

\label{idx:271}
有一个类似的概念——有时被视为 IaaS 的子集或扩展——是 \textbf{容器即服务 }(\textbf{CaaS})。在 CaaS 中，服务提供给你的不再是裸金属系统和虚拟机，而是容器与容器编排能力，你可以用它们构建自己的容器集群。DevOps 工程师可以直接与容器交互，按需对其进行配置与扩缩容。Google Cloud Platform 和 AWS 等都提供 CaaS 产品。
\end{myNotic}

\subsubsection{平台即服务模型}

\label{idx:272}
\label{idx:273}
如果基础设施本身无法满足你的需求，你可以改用 \textbf{Platform as a Service }(\textbf{PaaS}) 模型。在这种模型中，云服务提供商不仅管理基础设施（与 IaaS 一样），还管理操作系统、所需的中间件以及运行时——也就是你部署软件所依托的平台。

PaaS 解决方案通常还会为你提供应用版本管理、服务监控与发现、数据库管理、商业智能，甚至开发工具。

借助 PaaS，你在整条开发流水线中都能得到覆盖：从构建、测试到部署、更新和管理你的服务。不过，PaaS 解决方案比 IaaS 产品更贵。另一方面，由于整个平台都由服务商提供，你可以削减开发软件各部分的成本与时间，还能轻松地为分布在全球各地的开发团队提供相同的环境配置。

所有主流云服务提供商都有自己的产品，例如 Google App Engine 或 Azure App Service。也有一些独立产品，比如 Heroku。

\begin{myNotic}{注意}

\label{idx:274}
\label{idx:275}
除了更通用的 PaaS，还有 \textbf{Communications Platform as a Service }(\textbf{CPaaS})，在这种模式下，你会得到完整的通信后端，包括音频和视频，可以将它集成到自己的解决方案中。借助这项技术，你可以轻松提供支持视频的服务台，或者只需把实时聊天集成到应用中。
\end{myNotic}

\subsubsection{软件即服务模型}

\label{idx:276}
\label{idx:277}
有时你可能不想自己开发某个软件组件，只想直接使用现成的。 \textbf{Software as a Service }(\textbf{SaaS}) 基本上是为你提供一个托管的应用。使用 SaaS 时，你既不必操心基础设施，也不必操心构建在其上的平台，甚至不用操心软件本身。提供商负责整个软件栈的安装、运行、更新和维护，以及备份、许可和扩缩容。

\label{idx:278}
在 SaaS 模式下可以获得的软件种类相当丰富。例子从 Microsoft 365、Google Docs 这类办公套件，到 Slack 这类消息软件，再到 \textbf{customer relationship management }(\textbf{CRM}) 系统，甚至延伸到游戏类解决方案，比如云游戏服务，让你可以玩托管在云端的资源消耗型视频游戏。

通常，访问这类服务只需要一个浏览器，因此这是为员工提供远程办公能力的重要一步。

你可以创建自己的 SaaS 应用并提供给用户，既可以按自己的喜好部署，也可以通过 AWS Marketplace 之类的渠道发布。

另一方面，这种模式也有其取舍：一旦发生网络中断或服务中断，你就会失去对提供商软件以及自己数据的访问。

此外，缺乏直接控制权意味着数据安全与隐私由提供商负责，而不是你。而且，由于对软件栈的依赖，迁移到另一个提供商可能既复杂又昂贵，这就是所谓的供应商锁定。

\subsubsection{函数即服务模型与无服务器架构}

\label{idx:279}
\label{idx:280}
随着云原生技术的兴起，另一种日益流行的模式是 \textbf{Function as a Service }(\textbf{FaaS})，它是无服务器架构的一个子集。使用 FaaS，你会得到一个平台（与 PaaS 类似），可以在上面运行短生命周期的应用或函数。

使用 PaaS 时，通常你至少需要让服务的一个实例始终运行，而在 FaaS 中，你只在真正需要时才运行它们。这会增加处理请求的时间（以秒计），因为每次都需要启动函数。不过，其中一些请求可以被缓存，从而同时降低延迟和成本。说到成本，如果你长时间运行函数，FaaS 可能比 PaaS 贵得多，所以在设计系统时必须算好这笔账。

如果使用得当，FaaS 会把服务器从开发者眼前抽象掉，可以降低你的成本，还能提供更好的可扩展性，因为它可以基于事件而非资源来驱动。这种模式常用于运行预先调度或手动触发的任务、处理批量或流式数据，以及处理进来的不太紧急的请求。一些流行的 FaaS 提供商包括 AWS Lambda、Azure Functions、Genezio、Cloudflare Workers、Twilio Functions、IronFunctions、Netlify Functions、Vercel Functions 和 Google Cloud Functions。

尽管无服务器架构与云原生设计相关，但它本身也是一个热门话题。自 AWS Lambda、AWS Fargate、Google Cloud Run 和 Azure Functions 等 FaaS 或 CaaS 产品问世以来，它变得非常流行。

\label{idx:281}
无服务器在很大程度上是 Heroku 这类 PaaS 产品的演进。它把底层基础设施抽象掉，让开发者可以专注于应用，而不必纠结于基础设施选型。

\label{idx:282}
相比更早的 PaaS 解决方案，无服务器还有一个额外好处：你不必为没用到的东西付费。使用无服务器时，你通常不是为某个固定的服务级别付费，而是为已部署工作负载的实际执行时间付费。如果你只想每天运行某段代码一次，就不必为底层服务器支付月费。

虽然我们没有深入探讨无服务器的细节，但有一点很重要：它很少与 C++ 一起使用。就 FaaS 而言，目前只有 AWS Lambda 和 Google Cloud Functions 支持将 C++ 作为一种可选语言。由于容器与语言无关，你可以在 AWS Fargate、Azure Container Instances 或 Google Cloud Run 等 CaaS 产品中使用 C++ 应用和函数。

如果你想运行与 C++ 应用配合使用的非 C++ 辅助代码，无服务器函数对你可能仍然有用。维护任务和计划任务非常适合无服务器，它们通常不需要 C++ 二进制程序那样的性能或效率。

\label{idx:283}
C++ 也可以编译为 \textbf{WebAssembly }(\textbf{Wasm})，这是一种可移植的二进制代码格式，有助于在现代 Web 浏览器中运行高性能应用。这种格式已经超出浏览器范畴，进入了 AWS Lambda、Azure Functions、Fermyon Cloud、wasmCloud 和 Fastly Compute 等云环境。Fermyon 的 SpinKube 在 Kubernetes 上运行无服务器应用（Spin Apps）。此外，WasmEdge 可作为无服务器应用、嵌入式函数、微服务、智能合约和 IoT 设备的运行时。

现在我们已经介绍了本地部署和常见的云服务模式，接下来讨论人们在设计分布式系统时会犯的一些错误假设。

\subsection{避免分布式计算的谬误}

\label{idx:284}
\label{idx:285}
初涉分布式计算的人在开始设计这类系统时，往往会忘记或忽略这类系统的某些方面。尽管这些问题早在 90 年代就已被发现，但时至今日仍然具有现实意义。

下面各小节将逐一讨论这些谬误。

\subsubsection{网络是可靠的}

网络设备是为多年无故障运行而设计的。即便如此，仍有许多因素会导致丢包，例如停电、无线网络信号差、配置错误、有人绊到线缆，甚至动物咬断电线。举例来说，Google 不得不给其海底光缆加上凯夫拉保护层，因为光缆被鲨鱼啃咬（没错，是真的！）。你应该始终假设数据可能在网络的某个地方丢失。即便没有发生这种情况，线路另一端仍可能出现软件问题。

\label{idx:286}
为防范这类问题，务必制定自动重试失败网络请求的策略，并有办法处理常见的网络问题。重试时，尽量不要让对方过载，也不要重复提交同一笔事务。你可以使用本地应用队列或基于网络的消息队列来暂存失败的请求并重试发送。

诸如断路器模式这样的模式（我们将在本章后面介绍）也能帮上忙。哦对了，千万别只是无限等待，让每个失败的请求都白白占用资源。

\subsubsection{延迟为零}

无论是网络还是你运行的服务，即使在正常情况下，响应都需要一些时间。有时它们不得不花更长时间，尤其是在负载高于平均水平时。有时候，你的请求可能需要好几秒才能完成，而不是短短几毫秒。关于典型的延迟数字，可参考 \href{https://github.com/sirupsen/napkin-math\#numbers}{https://github.com/sirupsen/napkin-math\#numbers}。

\emph{图 4.2} 改编自 Jeff Dean 的著名演讲 \emph{Numbers Everyone Should Know}，它也直观地总结了在系统不同层次上观测到的典型延迟数字：

\myGraphic{0.920}{images/B18456_04_02.png}{图 4.2：Numbers Everyone Should Know，来自 Jeff Dean（\texthref{https://www.cs.cornell.edu/projects/ladis2009/talks/dean-keynote-ladis2009.pdf}）}

尽量把系统设计成不必等待过多细粒度的远程调用，因为每一次这样的调用都会增加总处理时间。即使在本地网络中，请求 1 条记录 10,000 次也会比请求 10,000 条记录 1 次慢得多。为降低网络延迟，可以考虑批量发送和处理请求。你也可以在等待小调用结果的同时执行其他处理任务，从而把这类调用的开销隐藏起来。

\label{idx:287}
\label{idx:288}
应对延迟的其他方式包括引入缓存、以发布者—订阅者模型主动推送数据而不是等待请求，或者将部署位置靠近客户，例如使用 \textbf{content delivery networks }(\textbf{CDNs})，即地理上分布式的代理服务器网络，通过降低网络延迟来帮助向用户分发内容。

\subsubsection{带宽是无限的}

在向架构中添加新服务时，务必留意它将占用多少流量。有时你可能希望通过压缩数据或引入限流策略来降低带宽占用。

这一谬误也与移动设备有关。如果信号较弱，网络往往会成为瓶颈。这意味着移动应用使用的数据量通常应保持在较低水平。使用 \emph{第 2 章}，\emph{架构风格} 中所述的 \emph{服务于前端的后端} 模式，往往有助于节省宝贵的带宽。

如果你的后端需要在某些组件之间传输大量数据，尽量确保这些组件彼此靠近：不要把它们运行在不同的数据中心。对于数据库来说，这通常归结为更好的复制。诸如 CQRS（本章后面会讨论）这样的模式也很有用。

\subsubsection{网络是安全的}

这是一个危险的谬误。一条链子的强度取决于它最薄弱的那一环，而遗憾的是，分布式系统中有很多这样的环节。以下是一些让这些环节更牢固的方法：

\begin{itemize}
\item 务必始终为基础设施的每一个组件、操作系统以及其他组件打上安全补丁。
\item 培训你的员工，并尽量保护系统免受人为因素影响；有时正是某个不守规矩的员工让系统遭到入侵。
\item 如果你的系统联网，它就会遭到攻击，而且某个时刻很可能会发生入侵。务必准备一份书面的保护方案，以便对此类事件作出响应。
\label{idx:289}
\item 你可能听说过 \textbf{纵深防御} 原则。它的核心是为系统的不同部分（你的基础设施、你的应用等）设置不同的检查，这样一旦发生入侵，其范围和造成的损害都会受到限制。
\item 使用防火墙、证书、加密以及恰当的认证。
\end{itemize}

关于安全的更多内容，请参阅 \emph{第 12 章}，\emph{代码与部署中的安全性}。

\subsubsection{拓扑不会改变}

这一谬误在微服务时代尤为明显。自动扩缩容以及管理基础设施时出现的 \emph{牲畜而非宠物} 方法（服务器好比可替换的牲畜，而不是独一无二的宠物）意味着拓扑会不断变化。这会影响延迟和带宽，也就是说，这一谬误带来的某些后果与前面讨论的谬误所导致的后果是一样的。

\emph{图 4.3} 展示了在持续变化的环境中可能出现的不同网络拓扑类型：

\myGraphic{0.730}{images/B18456_04_03.png}{图 4.3：网络拓扑类型}

\label{idx:290}
所幸，上述方法也附带了关于如何有效管理你的服务器 \emph{畜群} 的指导。依靠主机名和 DNS，而不是硬编码可能发生变化的 IP 地址，是朝着正确方向迈出的一步；而在 \emph{第 18 章}，\emph{云原生设计} 中介绍的服务发现，则是另一步。

第三步，也是更大的一步，是始终假设你的实例可能发生故障，并自动对此类情形作出反应。Netflix 的 \emph{Chaos Monkey} 工具也能帮你检验自身的准备程度，因为该工具会随机终止虚拟机和容器，从而模拟现实世界中的混乱，以验证服务是否具备故障韧性。

\subsubsection{只有一名管理员}

\label{idx:291}
关于分布式系统的知识，由于其自身的性质，往往也是分散的。这类系统及其基础设施的开发、配置、部署和管理由不同的人负责。不同的组件常常由不同的人升级，且不一定是同步进行的。还有所谓的 \emph{巴士系数}，简而言之，它是指某个关键项目成员被巴士撞到、从而导致维护项目或系统所必需的关键知识流失的风险因素。

我们该如何应对这一切？答案由几个部分组成。其中之一就是 DevOps 文化。通过促进开发与运维之间的紧密协作，人们共享关于系统的知识，从而降低巴士系数。此外，引入持续集成、持续交付和持续部署有助于升级项目并让它始终保持可用。持续集成与持续部署将在 \emph{第 11 章} 中介绍。

\label{idx:292}
你也可以尝试把系统设计成松耦合且向后兼容，这样升级某些组件时就不需要同时升级其他组件。一种简单的解耦方式是在它们之间引入消息传递，所以不妨考虑加一两个队列。这也有助于减少升级期间的停机时间。

\label{idx:293}
最后，尽量监控你的系统，并把日志集中收集到一个地方。系统去中心化并不意味着你现在需要手动查看十几台不同机器上的日志。\textbf{Elasticsearch, Logstash, Kibana }(\textbf{ELK}) 技术栈在这方面非常有价值。Grafana、Prometheus、Loki、Jaeger 和 OpenTelemetry Collector 也非常流行，尤其是在 Kubernetes 环境中。如果你在寻找比 Logstash 更轻量的方案，可以考虑 Fluentd 和 Filebeat，尤其是在处理容器时。这些工具将在 \emph{第 17 章}，\emph{可观测性} 和 \emph{第 18 章}，\emph{云原生设计} 中介绍。

\subsubsection{传输成本为零}

在规划项目及其预算时，考虑这一谬误非常重要。为分布式系统搭建和维护网络既费时又费钱，无论你是本地部署还是部署在云端——区别只在于你何时支付这笔成本。尽量估算设备成本、要传输的数据量（云服务提供商会对此收费）以及所需的人力。

如果你依赖数据压缩，要注意：虽然它降低了网络成本，却可能增加计算成本。一般来说，使用基于 gRPC 之类的二进制 API 会比基于 JSON 的 API 更便宜（也更快），而后者又比 XML 便宜，因为 JSON 和 XML 都是基于文本的冗长格式，而 gRPC 则大量依赖 \textbf{Protocol Buffers} (\textbf{Protobuf})，一种高效打包的二进制数据。如果你传输图像、音频或视频，就更必须估算这要花多少钱。

\subsubsection{网络是同质的}

即便你规划好了网络上要使用什么硬件、运行什么软件，最终也很容易出现一定程度的异构性。某些机器上的配置略有不同、需要集成的那个遗留系统使用了不同的通信协议，或者不同的手机向你的系统发送请求，都只是其中几个例子。另一个例子是通过云端额外的 worker 来扩展你的本地部署方案。

尽量限制所使用的通信协议和格式的数量，努力采用标准协议，因为实现互操作性很痛苦；同时避免供应商锁定，以确保你的系统在这种异构环境中仍能正常通信。

异构性也可能意味着韧性上的差异。这是因为同质网络的特征是使用相同的网络设备与协议（包括操作系统和应用），而异构网络则由这些组件的不同类型构成。可以尝试结合断路器模式与重试来处理这一点，这两者我们都会在本章后面讨论。

既然我们已经讨论了所有谬误，接下来讨论分布式架构的一个重要原则——\textbf{CAP 定理}——以及相关的 \textbf{最终一致性} 概念。

\subsection{CAP 定理与最终一致性}

\label{idx:294}
\label{idx:295}
要设计出成功跨越多个节点的系统，你需要了解并运用某些原则。其中之一就是 CAP 定理。它涉及设计分布式系统时最重要的选择之一，其名称源自分布式系统可以具备的三个属性。它们如下：

\begin{itemize}
\item \textbf{一致性}：每次读取都会让你得到最近一次写入之后的数据（或者返回一个错误）
\item \textbf{可用性}：每个请求都会得到非错误的响应（但不保证你拿到的是最新数据）
\item \textbf{分区容错性}：即使两个节点之间发生网络故障，整个系统仍会继续工作
\end{itemize}

本质上，该定理指出，对于分布式系统，你最多只能从这三个属性中选取两个：

\myGraphic{0.619}{images/B18456_04_04.png}{图 4.4：CAP 定理，又称 Brewer 定理}

只要系统正常运行，看起来这三个属性都能得到满足。然而，正如我们从那些谬误中了解到的，网络并不可靠，因此分区终会发生。在这种情况下，分布式系统仍应正常运行。这意味着该定理迫使你在提供分区容错性与一致性（即 CP）还是分区容错性与可用性（即 AP）之间做出选择。通常后者是更好的选择。如果你想选择 CA，就必须彻底移除其他节点，只留下一个单节点系统。

如果在发生分区时你决定提供一致性，就必须要么返回错误，要么在等待数据达到一致的过程中承受超时风险。如果你选择可用性而非一致性，就要冒返回过期数据的风险——最新的写入可能无法跨分区传播。

这两种取舍适用于不同的需求。如果你的系统需要原子性的读写（\emph{一次一笔}），例如因为客户可能会损失钱财，那就选择 CP。如果你的系统必须在分区下继续运行，或者你可以接受最终一致性，那就选择 AP。

好，但什么是最终一致性？下面我们来讨论一致性的不同级别，以便理解它。

\label{idx:296}
\label{idx:297}
在提供强一致性的系统中，每次写入都是同步传播的。这意味着所有读取总是能看到最新的写入，代价则是更高的延迟或更低的可用性。这正是 \textbf{relational database management systems }(\textbf{RDBMS}) 所提供的类型（基于 \textbf{ACID} 保证，即 \textbf{原子性、一致性、隔离性、持久性}），最适合需要事务的系统。

\label{idx:298}
另一方面，在提供最终一致性的系统中，你只保证一次写入之后，读取最终会看到这一变化。通常，\emph{最终} 意味着在几毫秒之内。这是由于这类系统中数据复制的异步特性，与上一段所说的同步传播正好相反。这里提供的不再是 ACID 保证（例如使用 RDBMS 时那样），而是 \textbf{BASE} 语义（意为 \textbf{基本可用、软状态、最终一致性}），通常由 NoSQL 数据库提供。

一个系统要做到异步且最终一致（AP 系统往往如此），就需要有解决状态冲突的办法。一种常见做法是在实例之间交换更新，并在同步更新时选择把最先或最后的写入作为被接受的版本。

\label{idx:299}
CAP 定理曾因过于简化重要概念而受到批评，这促成了 BASE 的提出以及作为 CAP 扩展的 \textbf{PACELC 定理} 的出现。PACELC 定理的表述为 \emph{IF P -> (A or C), ELSE (L or C)}。也就是说，如果存在网络分区，系统就要在一致性与可用性之间选择；否则，系统要在延迟与一致性之间选择。该定理刻画了分布式系统中在时间延迟与数据一致性之间寻求权衡的必要性，并且比 CAP 更准确地划分 NoSQL 数据库。

\myGraphic{0.664}{images/B18456_04_05.png}{图 4.5：PACELC 定理}

下面我们来讨论两个有助于实现最终一致性的相关模式。

\subsubsection{Saga 模式与补偿事务模式}

当你需要执行分布式事务时，saga 模式就很有用。

\label{idx:300}
在微服务时代之前，如果只有一台主机和一个数据库，你可以依赖数据库引擎为你完成事务。如果一台主机上有多个数据库，你可以使用\textbf{两阶段提交}（\textbf{2PC}）来做到这一点。使用 2PC 时，你会有一个协调者，它先通知所有数据库进行准备，等它们都报告准备就绪后，再通知它们全部提交事务。

而现在，每个微服务通常都有自己的数据库（若想获得可扩展性，也理应如此），并且它们分散在整个基础设施中，你再也无法依赖简单的事务和 2PC 了（失去这种能力往往意味着你不再需要 RDBMS，因为 NoSQL 数据库可以快得多）。

相反，你可以使用 saga 模式。下面用一个例子来说明。

假设你想创建一个在线仓库，它能追踪库存量并支持信用卡支付。要处理一笔订单，你至少需要三个服务：一个处理订单，一个预留库存，还有一个向信用卡扣款。

\label{idx:301}
\label{idx:302}
\label{idx:303}
\label{idx:304}
如今，saga 模式有两种实现方式：基于协同（choreography）的（也称基于事件的）和基于编排（orchestration）的（也称基于命令的）。

\paragraph{基于协同的 saga}

\label{idx:305}
在协同方式中，服务之间交换事件，没有中心化的控制器。在这种情况下，saga 的第一部分是由订单处理服务向库存服务发送一个事件。库存服务完成自己的部分后，再向支付服务发送另一个事件。支付服务接着又向订单服务发回一个事件。这样就完成了事务（即 saga），订单现在可以顺利发货了。

如果订单服务想要追踪事务的状态，它只需同样监听所有这些事件即可。

当然，有时订单无法完成，就需要进行回滚。在这种情况下，saga 的每一步都需要分别且谨慎地回滚，因为其他事务可能正在并行运行，例如修改库存状态。这类回滚称为补偿事务（compensating transaction），用于撤销操作前几个阶段已完成的所有工作。微服务采用这种方式，是因为把单体数据库拆分成多个数据库后，ACID 事务已不可能实现。

\myGraphic{0.891}{images/B18456_04_06.png}{图 4.6：协同}

这种实现 saga 模式的方式相当直接，但如果涉及的服务之间存在很多依赖关系，那么采用\emph{编排}方式可能更好。

\paragraph{基于编排的 saga}

\label{idx:306}
在编排方式中，服务依赖一个中心化的控制器（编排器），由它根据事件告诉各个服务该执行哪些操作。在这种情况下，我们可以使用消息代理来处理服务与编排器之间的通信，由编排器来协调 saga。订单服务会向编排器发送一个请求，编排器随后向库存服务和支付服务发送命令。这两者各自完成自己的部分后，再通过代理上可用的回复通道把应答发回编排器。

在这种场景下，编排器拥有协调事务所需的全部逻辑，而各个服务本身无需知道还有哪些其他服务参与了该 saga。

如果编排器收到某个服务失败的消息（例如信用卡已过期），它就需要启动回滚。在我们的例子中，它会再次借助代理，向特定服务发送相应的回滚命令。

\myGraphic{0.648}{images/B18456_04_07.png}{图 4.7：编排}

\label{idx:307}
好了，关于最终一致性就讲到这里。下面我们来转向与可用性相关的其他话题。

\section{让系统具备容错性和可用性}

\label{idx:308}
可用性和容错性是对每种架构都至少有一定重要性的软件质量属性。如果系统无法访问，创建软件系统又有什么意义呢？本节我们将了解这两个术语的确切含义，以及一些将它们融入你解决方案的技术。

\subsection{计算系统的可用性}

\label{idx:309}
可用性是指系统处于运行、可用且可访问状态的时间百分比。导致系统无法响应的崩溃、网络故障或极高负载（例如来自 DDoS 攻击的负载）都会影响其可用性。

通常，追求尽可能高的可用性水平是个好主意。你可能会遇到\emph{数 9} 这个说法，因为可用性常常表示为 99\%（两个 9）、99.9\%（三个 9）等等。每多一个 9 都更难获得，所以做承诺时要谨慎。请看下表，了解不同可用性水平下每月可接受的停机时间：

\begin{center}
\begin{tabularx}{\textwidth}{XX}
\toprule
\textbf{停机时间/月} & \textbf{正常运行时间} \\
\midrule
7 小时 18 分钟 & 99\%（两个 9） \\
43 分钟 48 秒 & 99.9\%（三个 9） \\
4 分钟 22.8 秒 & 99.99\%（四个 9） \\
26.28 秒 & 99.999\%（五个 9） \\
2.628 秒 & 99.9999\%（六个 9） \\
262.8 ms & 99.99999\%（七个 9） \\
26.28 ms & 99.999999\%（八个 9） \\
.628 ms & 99.9999999\%（九个 9） \\
\bottomrule
\end{tabularx}
\end{center}

表 4.1：不同可用性水平对应的每月停机时间

\label{idx:310}
云应用的一个常见做法是提供\textbf{服务级别协议}（\textbf{SLA}），其中规定在给定时间段（例如一年）内允许发生多少停机时间。你的云服务的 SLA 在很大程度上取决于你所依赖的云服务的 SLA。

要计算两个需要协作的服务的复合可用性，只需将它们的正常运行时间相乘。也就是说，如果你有两个可用性为 99.99\% 的服务，它们的复合可用性就是 \emph{99.99\% * 99.99\% = 99.98\%}。要计算冗余服务（例如两个相互独立的区域）的可用性，则应将它们的不可用时间相乘。例如，如果两个区域的可用性都是 99.99\%，它们的总不可用时间就是 \emph{(100\% – 99.99\%) * (100\% – 99.99\%) = 0.01\% * 0.01\% = 0.0001\%}，因此它们的复合可用性为 99.9999\%。

\label{idx:311}
遗憾的是，100\% 的可用性是不可能实现的。故障总会时不时发生，所以下面来学习如何让你的系统容忍它们。

\subsection{构建容错系统}

\label{idx:312}
容错性是指系统检测此类故障并优雅地处理它们的能力。你的云服务必须具备韧性，这一点至关重要，因为由于云的特性，许多意想不到的情况可能突然出问题。良好的容错性能有助于提升服务的可用性。

不同类型的问题需要不同的处理方式：从预防，到检测，再到尽量减小影响。我们先从避免出现单点故障的常见做法说起。

\subsubsection{冗余模式}

\label{idx:313}
\label{idx:314}
最基本的预防手段之一是引入冗余。就像你可以为汽车准备一个备胎一样，你也可以准备一个备用服务，在主服务器宕机时接管工作。这种接管也称作\textbf{故障转移}（failover）。

备用服务器如何知道何时该接管呢？实现这一点的一种方式是使用\emph{检测故障}一节中介绍的心跳机制。

\label{idx:315}
\label{idx:316}
为了让切换更快，你可以把发往主服务器的所有消息也发送到备用服务器。这称为\textbf{热备}（hot standby），与之相对的是从零初始化的冷备。这种情况下，一个不错的做法是滞后一条消息：这样如果某条\emph{被投毒}的消息杀死了主服务器，备用服务器只需拒绝它即可。

\label{idx:317}
上述机制称为\textbf{主动—被动}（或\textbf{主—从}）故障转移，因为备用服务器不处理进入的流量。如果它也处理，我们就有了\textbf{主动—主动}（或\textbf{主—主}）故障转移。有关主动—主动架构的更多内容，请参阅\emph{延伸阅读}一节。

\myGraphic{0.920}{images/B18456_04_08.png}{图 4.8：主动—主动与主动—被动}

\label{idx:318}
请确保在发生故障转移时不会丢失任何数据。使用带持久化存储的消息队列来重新提交消息进行处理，可能有助于做到这一点。

\paragraph{领导者选举}

\label{idx:319}
\label{idx:320}
\label{idx:321}
让各服务器明确谁是谁也很重要——如果不止一个服务器开始表现得像主实例，就会导致失步和数据丢失，因为其中一个节点是分布在其他节点（计算机）上的任务的组织者。选择主服务器的过程称为领导者选举模式。实现它的算法有几种，例如引入第三方仲裁者、通过竞争获取某个共享资源的独占所有权、选择秩最低的实例，或使用 bully 选举（霸道选举）、令牌环选举等算法。

\myGraphic{0.542}{images/B18456_04_09.png}{图 4.9：领导者选举}

领导者选举也是下一个相关概念——达成共识——的重要组成部分。

\paragraph{共识}

\label{idx:322}
如果你希望系统在网络分区发生或服务的某些实例出现故障时仍能运行，就需要一种让各实例达成共识的方式。它们必须就提交哪些值、以及（通常还包括）以何种顺序提交达成一致。一种简单的方法是让每个实例对正确的状态投票。然而，在某些情况下，这不足以正确达成共识，甚至根本无法达成。另一种方法是选出一个领导者，让它传播自己的值。由于手工实现这类算法并不容易，我们建议使用 Paxos、Raft 等久经业界验证的流行共识协议。后者正日益流行，因为它更简单、更易于理解。

下面我们来讨论另一种预防系统故障的方法。

\subsubsection{复制模式}

\label{idx:323}
\label{idx:324}
这种方式在数据库中尤其流行，也有助于数据库的扩展。复制是指并行运行服务的多个实例，它们持有相同的数据副本，并共同处理进入的流量。

\begin{myNotic}{注意}

不要把复制与分片混为一谈。分片不需要任何数据冗余，但在大规模场景下往往能带来出色的性能。如果你使用的是 Postgres，我们建议你试试 Citus（\href{https://www.citusdata.com/}{https://www.citusdata.com/}）、CockroachDB（\href{https://www.cockroachlabs.com/}{https://www.cockroachlabs.com/}）或 YugabyteDB（\href{https://www.yugabyte.com/}{https://www.yugabyte.com/}）。
\end{myNotic}

就数据库而言，复制有两种方式。

\paragraph{主从复制}

\label{idx:325}
在这种场景下，所有服务器都能执行只读操作，但只有一个主服务器还可以执行写入。数据从主服务器经由从服务器复制，拓扑结构为一对多或树形。如果主服务器发生故障，系统仍能以只读模式运行，直到该故障被修复。

\myGraphic{0.772}{images/B18456_04_10.png}{图 4.10：主从复制}

\paragraph{多主复制}

\label{idx:326}
你也可以构建一个拥有多个主服务器的系统。如果有两个服务器，就构成主主复制方案。如果其中一个服务器宕机，其他服务器仍能正常运行。不过，此时你要么需要同步写入，要么就得提供更宽松的一致性保证。此外，你还需要提供一个路由器。

\myGraphic{0.920}{images/B18456_04_11.png}{图 4.11：多主复制}

\label{idx:327}
这类复制的例子包括 Microsoft 的 Active Directory、OpenLDAP 和 Apache CouchDB。

下面我们来讨论两种预防由过高负载引发故障的方法。

\subsubsection{基于队列的负载均衡模式}

\label{idx:328}
\label{idx:329}
这种策略旨在降低系统负载突然激增所造成的影响。用大量请求淹没某个服务会导致性能下降、可靠性降低，甚至丢失有效请求。队列又一次派上了用场，解了燃眉之急。

要实现这一模式，我们只需引入一个队列，让进入的请求被异步地加入其中。你可以使用 Amazon 的 SQS、Azure 的 Service Bus、Apache Kafka、ZeroMQ、RabbitMQ、EMQX、HiveMQ 或其他队列来做到这一点。

这样一来，进入的请求不再出现峰值，负载被平均化了。我们的服务可以从上述队列中取出请求并加以处理，甚至根本不知道负载曾经升高。就这么简单。

如果你的队列性能出色，而且任务可以并行化，那么这一模式还有一个附带的好处：更好的可扩展性。

此外，如果你的服务不可用，请求仍会被加入队列，等待该服务恢复后再处理，因此这也许有助于提升可用性。

如果请求到来得并不频繁，可以考虑把服务实现为无服务器函数，仅在队列中有待处理项时才运行，以节省成本。

请记住，使用这一模式时，由于增加了一个队列，整体延迟会上升。Apache Kafka 和 ZeroMQ 应当能带来较低的延迟，但如果这一点让你无法接受，还有另一种应对负载增加的办法。

\subsubsection{背压模式}

\label{idx:330}
\label{idx:331}
如果负载持续居高不下，你的任务很可能超出你的处理能力。当请求再也无法全部放入内存时，这会导致缓存未命中和内存交换，还会造成请求被丢弃以及其他种种糟糕情况。如果你预计会有高负载，施加背压也许是一种很好的应对方式。

本质上，背压意味着我们不是随着每个进入的请求给服务施加更多压力，而是把压力推回给调用方，让它去处理这种局面。实现这一点有几种不同的方式。

例如，我们可以阻塞接收网络数据包的线程。调用方就会发现自己无法把请求推送给我们的服务——我们反而把压力向上游推了回去。

\label{idx:332}
\label{idx:333}
另一种方式是识别出负载过高，直接返回一个错误码，例如 \texttt{503}。你可以这样设计架构，让另一个服务替你完成这件事。Envoy Proxy（\href{https://envoyproxy.io}{https://envoyproxy.io}）就是这样一个服务，它在许多其他场合也很好用。

Envoy 是一款用 C++ 编写的高性能代理服务器，专为微服务架构而设计。它管理微服务之间的网络通信，充当客户端与服务之间的中间层。它提供高效的路由、负载均衡、增强的安全性、监控和请求追踪。特别是，Istio 是一个服务网格，它借助 Envoy 来中介服务的所有入站和出站流量。

\label{idx:334}
Envoy 可以基于预定义的配额施加背压，因此你的服务实际上永远不会过载。它还可以测量处理请求所需的时间，只有在超过某个阈值时才施加背压。还有许多其他情况会返回各种错误码。但愿调用方对压力回到自己身上时该怎么做已有所准备。

\begin{myNotic}{注意}

请注意，本章提到 Envoy 并非是在指定你必须使用它，而只是用它来演示所讨论的模式在现实系统中是如何应用的。使用其他工具同样可以达到类似的效果。
\end{myNotic}

既然已经知道如何预防故障，下面来学习故障发生后如何检测它们。

\subsection{检测故障}

\label{idx:335}
恰当而快速的故障检测能为你省去很多麻烦，往往也能省钱。针对不同需求，检测故障的方式多种多样。下面我们来介绍其中的几种。

\subsubsection{边车设计模式}

\label{idx:336}
既然我们一直在讨论 Envoy，不妨提一句：它就是边车设计模式的一个例子。这一模式的用途远不止错误预防和检测，Envoy 就是一个很好的例证。

\label{idx:337}
总的来说，边车让你无需编写额外代码就能为服务增添许多能力。就像物理的边车可以挂在摩托车上一样，软件边车也可以挂接到你的服务上，二者都能扩展所提供的功能。

\label{idx:338}
边车在检测故障方面能起什么作用呢？首先，它能提供健康检查能力。就被动健康检查而言，Envoy 可以检测出服务集群中是否有某个实例开始表现异常。这称为离群点检测。Envoy 可以查找连续的 \texttt{5XX} 错误码、网关故障等。除了检测出这类有问题的实例，它还能将其剔除，从而让整个集群保持健康。

Envoy 还提供主动健康检查，也就是说，它可以直接探测服务本身，而不只是观察服务对进入流量的反应。

\label{idx:339}
\label{idx:340}
在本章中，我们还会介绍边车模式（尤其是 Envoy）的其他一些用途。下面我们来讨论下一种故障检测机制。

\subsubsection{心跳机制模式}

\label{idx:341}
最常见的故障检测方式之一是通过心跳机制。心跳是在两个服务之间按固定时间间隔（通常为几秒）发送的信号或消息。

\label{idx:342}
如果连续丢失了几个心跳，接收方服务就可以认定发送方服务\emph{已死亡}。就前面几节提到的那个主备服务对而言，这会触发一次故障转移。

实现心跳机制时，务必确保它可靠。误报会很麻烦，因为服务可能会混淆，例如搞不清谁应该成为新的主服务器。一个不错的做法是专门为心跳提供一个独立端点，这样它就不容易受到常规端点流量的影响，因为服务必须尽快响应心跳请求。

\subsubsection{漏桶计数器模式}

\label{idx:343}
\label{idx:344}
另一种检测故障的方法是引入所谓的漏桶计数器。每发生一次错误，计数器就递增；当达到某个阈值（桶满了）后，就发出故障信号并进行处理。计数器会按固定时间间隔递减（因此称为漏桶）。这样一来，只有在短时间内发生大量错误时，才会被判定为故障。

如果在你的场景中偶尔出现错误是正常的（例如涉及网络通信时），这一模式会很有用。

既然已经知道如何检测故障，下面来学习故障发生后该怎么做。

\subsection{尽量减小故障的影响}

\label{idx:345}
检测一个正在发生的故障需要时间，解决它则要耗费更多这种宝贵资源。正因如此，你应当努力把故障的影响降到最低。以下有几种行之有效的方法。

\subsubsection{重试调用（重试模式）}

\label{idx:346}
当应用调用另一个服务时，调用有时会失败。针对这种情况，最简单的补救办法就是重试该调用。如果故障是暂时性的而你没有重试，该故障很可能会在系统中传播，造成本不该有的更大破坏。实现一种自动重试此类调用的方式，能为你省去很多麻烦。

还记得我们的边车代理 Envoy 吗？原来它可以替你执行自动重试，让你无需对源代码做任何改动。

例如，看看下面这个可添加到 Envoy 路由中的重试策略示例配置：

\begin{codebox}[yaml]
retry_policy:
  retry_on: "5xx"
  num_retries: 3
  per_try_timeout: 2s
\end{codebox}

\label{idx:347}
这会让 Envoy 在调用返回错误（例如 \texttt{503} HTTP 状态码或映射到 \texttt{5XX} 的 gRPC 错误）时重试该调用。总共会重试三次，每次若未在 2 秒内完成就算失败。

\subsubsection{避免级联故障}

我们提到过，如果不重试，错误就会传播，在整个系统中引发一连串故障。下面再介绍一些防止这种情况发生的方法。

\paragraph{断路器模式}

\label{idx:348}
\label{idx:349}
断路器模式是应对这一问题的非常有用的工具。它让我们能够迅速发现某个服务无法处理请求，从而可以短路对该服务的调用。这既可以发生在靠近被调用方的一侧（Envoy 提供了这种能力），也可以发生在调用方一侧（额外好处是能缩短调用耗时）。就 Envoy 而言，只需将以下内容添加到配置中即可：

\begin{codebox}[yaml]
circuit_breakers:
  thresholds:
    - priority: DEFAULT
      max_connections: 1000
      max_requests: 1000
      max_pending_requests: 1000
\end{codebox}

在这两种情况下，对服务发起调用所造成的负载都可能下降，这在某些情况下有助于服务恢复正常运行。

如何在调用方一侧实现断路器呢？在发起若干次调用之后，比如说你的漏桶溢出了，你只需在指定的一段时间内停止发起新的调用（例如等到桶不再溢出为止）。简单而有效。

\paragraph{舱壁模式}

\label{idx:350}
另一种限制故障扩散的方法直接取自造船业。造船时，你通常不希望在船体破洞后整艘船都灌满水。为限制这类破洞造成的损害，你会把船体分隔成若干舱壁，每个舱壁都易于隔离。这样一来，只有破损的那个舱壁会进水。

\label{idx:351}
同样的原理也适用于在软件架构中限制故障影响。你可以把实例划分成若干组，也可以把它们使用的资源划分到组中。设置配额同样可以被视为这一模式的一个例子。

可以为不同的用户群体创建相互独立的舱壁，当你需要对它们进行优先级排序，或为关键消费者提供不同级别的服务时，这会很有用。

\subsubsection{Geodes 模式}

\label{idx:352}
\label{idx:353}
我们要讨论的最后一种降低故障影响的方法称为\textbf{geodes}。这个名字来自「地理节点」（geographical nodes）。当你的服务部署在多个区域时可以使用它。

\label{idx:354}
如果某个区域发生故障，你只需把流量重定向到其他未受影响的区域即可。当然，这会让延迟比调用同一数据中心内其他节点时高得多，但通常来说，把不太关键的用户重定向到远程区域，要比让他们的调用直接失败好得多。

现在你已经知道如何通过系统架构来提供可用性和容错，下面我们来讨论如何把系统的各个组件集成到一起。

\section{集成你的系统}

\label{idx:355}
分布式系统并不只是一堆彼此孤立、对外界一无所知地运行着的应用实例的集合。这些实例不断地相互通信，必须被妥善集成才能发挥最大价值。

关于集成这一主题已有许多论述，因此本节我们只展示寥寥几种模式，用于高效地集成全新的系统，以及需要与现有部分（往往是遗留部分）共存的新系统组件。

为了避免让本章自己就写成一本完整的书，我们先以一份延伸阅读推荐来开启本节。如果你对集成模式感兴趣，尤其是以消息传递为重点的集成模式，那么 Gregor Hohpe 和 Bobby Woolf 的《企业集成模式》（Enterprise Integration Patterns）一书就是你的必读之作。

下面我们简要看看该书涵盖的两种模式。

\subsection{管道与过滤器模式}

\label{idx:356}
\label{idx:357}
\label{idx:358}
\label{idx:359}
我们要讨论的第一种集成模式称为管道与过滤器。它的目的是把一项庞大的处理任务分解成一系列更小、相互独立的任务（称为过滤器），然后你可以把它们连接起来（使用管道，例如消息队列）。这种方式能带给你可扩展性、性能和可复用性。

\myGraphic{0.742}{images/B18456_04_12.png}{图 4.12：管道与过滤器}

假设你需要接收并处理一笔进入的订单。你可以把它放在一个大模块里完成，这样就不需要额外的通信，但这样一个模块的不同功能会难以测试，也更难做好扩展。

\label{idx:360}
\label{idx:361}
相反，你可以把订单处理拆分成若干独立的步骤，每一步由一个单独的组件负责：一个负责解码，一个负责校验，另一个负责订单的实际处理，再有一个负责把订单存储到某处。采用这种方式，你如今可以独立执行其中的每一步，在需要时轻松替换或停用它们，并将它们复用来处理不同类型的输入消息。

如果你想同时处理多个订单，也可以把处理过程流水线化：当一个线程校验某条消息时，另一个线程解码下一条消息，依此类推。

其缺点在于，你需要用同步队列来充当管道，这会带来一些额外开销。

若要扩展处理过程中的某一步，你可能需要把这种模式与我们接下来要介绍的另一种模式结合使用。

\subsection{竞争消费者模式}

\label{idx:362}
\label{idx:363}
竞争消费者的思路很简单：你有一个输入队列（或消息通道），以及若干个消费者实例，它们并发地从队列中取出并处理条目。每个消费者都能处理消息，因此它们会相互竞争，看谁能成为接收者。

\myGraphic{0.846}{images/B18456_04_13.png}{图 4.13：竞争消费者}

这样一来，你便获得了可扩展性、免费的负载均衡以及韧性。再加上这个队列，你同时还具备了基于队列的负载均衡模式。

如果你需要削减某个请求的延迟，或者只是希望提交到队列中的某个特定任务能更紧急地执行，这种模式可以轻而易举地与优先队列结合使用。

\begin{myNotic}{注意}

如果顺序很重要，这种模式用起来可能会比较棘手。消费者接收消息并完成处理的顺序可能会有所不同，所以要确保这不会影响你的系统，或者你要有办法稍后对结果重新排序。如果你需要按顺序处理消息，可能就无法使用这种模式。
\end{myNotic}

\label{idx:364}
\label{idx:365}
下面再看几种模式，这次是为了帮助我们把新系统与现有系统集成起来。

\subsection{从遗留系统迁移}

\label{idx:366}
从零开始开发一个系统可以是一段幸福的体验。做的是开发而不是维护，还有机会使用最前沿的技术栈——有什么理由不喜欢呢？可惜的是，一旦开始与现有的遗留系统集成，这种幸福往往就结束了。不过幸运的是，有一些方法可以减轻这种痛苦。

\subsubsection{防腐层模式}

\label{idx:367}
\label{idx:368}
引入防腐层可以帮助你的解决方案无痛地与语义不同的遗留系统集成。这个额外的层负责两侧之间的通信。

这样的组件让你的解决方案在设计上拥有更大的灵活性，无需在技术栈或架构决策上做出妥协。要做到这一点，只需对遗留系统做极少的改动（如果遗留系统不需要调用新系统，则完全无需改动）。

\myGraphic{0.920}{images/B18456_04_14.png}{图 4.14：防腐层}

例如，如果你的解决方案基于微服务，遗留系统只需与防腐层通信，而不必定位并直接访问每一个微服务。任何转换（例如由于协议版本过时）也在这个额外的层中完成。

\label{idx:369}
\label{idx:370}
请记住，添加这样的层会引入延迟，并且必须满足你解决方案的质量属性要求，例如可扩展性。

\subsubsection{绞杀者模式}

\label{idx:371}
绞杀者模式允许从遗留系统逐步迁移到新系统。我们刚看过的防腐层有助于两个系统之间的通信，而绞杀者模式则用于由两者共同向外部世界提供服务。

\label{idx:372}
在迁移过程的初期，绞杀者外观（一个代理或 API 网关）会把大部分请求路由到遗留系统。在迁移过程中，越来越多的调用可以被转发到新系统，同时越来越彻底地 \emph{绞杀} 遗留系统，限制它所提供的功能。作为迁移的最后一步，绞杀者连同遗留系统可以一并退役——此时新系统将提供全部功能：

\myGraphic{0.920}{images/B18456_04_15.png}{图 4.15：对单体应用的绞杀。迁移之后，绞杀者仍可用作遗留请求的入口点或适配器}

对于小型系统来说，这种模式可能过于复杂；如果数据存储需要共享，或者系统采用事件溯源，它也可能变得棘手。把它加入你的解决方案时，一定要规划好如何达到合适的性能和可扩展性。

说到这两个属性，下面我们来讨论一些有助于实现它们的方法。

\section{实现大规模下的性能}

在设计 C++ 应用时，性能通常是一个关键因素。虽然使用这门语言本身在单个应用范围内已经能带来很大帮助，但恰当的高层设计对于在规模化时实现最优的延迟和吞吐量同样不可或缺。下面我们来讨论几个有助于此的关键模式与方面。

\subsection{CQRS 与事件溯源模式}

\label{idx:373}
扩展计算有许多方法，但扩展数据访问可能比较棘手。然而，当用户规模增长时，这往往又是必需的。\textbf{命令查询职责分离 }（\textbf{CQRS}）就是一种能在这方面提供帮助的模式。

\subsubsection{CQRS}

\label{idx:374}
在传统的\textbf{创建、读取、更新、删除}（\textbf{CRUD}）系统中，读和写都使用同一个数据模型，数据也以相同的方式流动。标题中所说的分离，基本上就是指把查询（读）和命令（写）用两种各自独立的方式来处理。

\myGraphic{0.781}{images/B18456_04_16.png}{图 4.16：CRUD 服务}

\label{idx:375}
许多应用的读写比例严重失衡——在典型应用中，从数据库读取的次数通常远多于更新它的次数。这意味着让读取尽可能快可以带来更好的性能：现在读和写可以分别优化和扩展。CQRS 可用于支持这一点。除此之外，当大量写入相互竞争、需要维护所有写入的记录，或者你的某部分 API 用户应当只有只读访问权限时，引入 CQRS 也会有所帮助。

\label{idx:376}
为读和写分别使用不同的模型，可以让不同的团队分别处理两侧的工作。负责读取一侧的开发者不需要深入理解领域（而这是正确执行更新所必需的）。当他们发起请求时，只需一次简单的调用就能从薄薄的读取层获得一个\textbf{数据传输对象 }（\textbf{DTO}），而不必经过领域模型。

如果你还不清楚 DTO 是什么，可以想想从数据库返回商品数据的场景。如果调用方想要一个商品列表，你可以提供一个只包含商品名称和缩略图的 \texttt{Item\allowbreak{}Overview} 对象。另一方面，如果他们想要某个特定门店的商品，你也可以提供一个 \texttt{StoreItem} 对象，其中包含名称、更多图片、描述和价格。\texttt{Item\allowbreak{}Overview} 和 \texttt{StoreItem} 都是 DTO，它们从数据库里同一批 \texttt{Item} 对象中获取数据。

读取层既可以位于写入所用的数据存储之上，也可以是一个通过事件来更新的不同数据存储，如下图所示：

\myGraphic{0.920}{images/B18456_04_17.png}{图 4.17：结合事件溯源的 CQRS；图中的领域存储是一个事件源}

\label{idx:377}
\label{idx:378}
采用 \emph{图 4.17} 所示的方法，你可以按需创建任意多种命令，每种命令都有自己的处理器。通常，命令是异步的，不会向调用方返回任何值。每个处理器都使用领域对象并持久化所做的更改。完成之后，会发布事件，事件处理器可以利用这些事件来更新读操作所用的存储。继续上一个例子，商品数据查询会从由 \texttt{ItemAdded} 或 \texttt{Item\allowbreak{}Price\allowbreak{}Changed} 等事件更新的数据库中获取信息，而这些事件可由 \texttt{AddItem} 或 \texttt{Modify\allowbreak{}Item} 等命令触发。

使用 CQRS 让我们可以为读操作和写操作使用不同的数据模型。例如，你可以创建存储过程和物化视图来加速读取。为领域存储和读存储使用不同类型的存储（SQL 或 NoSQL）也可能是有益的，因为这些存储类型可能分别更适合更新数据或读取数据。一种高效的数据持久化方式是使用 Apache Cassandra 集群，而使用 Elasticsearch 则是快速搜索已存储数据的好方法。

\label{idx:379}
除了前面提到的优点，CQRS 也有它的缺点。由于它引入了复杂性，通常并不适合小型或要求不高的架构。通常有用的做法是只把它应用到系统中能带来最大收益的部分。你还应该注意到，在领域存储之后再更新读存储，意味着我们现在拥有的是最终一致性，而不是强一致性。

\subsubsection{命令查询分离}

\label{idx:380}
CQRS 基于一个很早就在 Eiffel 编程语言（也是引入契约概念的那门语言）中提出的简单思想：\textbf{命令查询分离 }（\textbf{CQS}）。CQS 是一条原则，主张把 API 调用分离成——嗯——命令和查询，就像 CQRS 一样，但不限规模。它与面向对象编程以及一般意义上的命令式编程配合得非常好。

如果你的函数名以 \texttt{has}、\texttt{is}、\texttt{can} 或类似的词开头，那么它就应该只是一个查询，不修改底层状态，也不产生任何副作用。这样做带来两大好处：

\begin{itemize}
\item \textbf{更容易推理代码}：这类函数在语义上显然只是 \emph{读取}，而绝不是 \emph{写入}。这在调试时能让查找状态变更变得容易得多。
\item \textbf{更少的海森堡 bug}：如果你曾经调试过一个只在发布版中显现而在调试版中不显现的错误（或者反过来），那你就是在和海森堡 bug 打交道——这类 bug 会在你调试软件时消失或改变行为。这很少是什么令人愉快的事。许多此类错误可能由修改状态的 assert 调用引起。遵循 CQS 可以消除这类 bug。
\end{itemize}

与 assert 类似，如果你想使用契约（前置条件和后置条件），极其重要的是在其中只使用查询。否则，禁用某些契约检查同样可能导致海森堡 bug，更别提这会多么违反直觉了。

下面我们再谈几句事件溯源。

\subsubsection{事件溯源}

\label{idx:381}
如 \emph{第 2 章} \emph{架构风格}中所述，事件溯源意味着：不必总是存储应用程序的整个状态（还可能要在更新时处理冲突），你只需存储应用程序状态所发生的变更。使用事件溯源可以消除并发更新，并让所有相关方都能对自己的状态进行渐进式变更，从而提升应用的性能。保存已执行操作的历史记录（例如市场交易）可以让调试（通过稍后重放它们）和审计更加容易。这还带来了更大的灵活性和可扩展性。引入事件溯源后，一些领域模型可以通过降低复杂度而变得简单得多。

事件溯源的一个代价是最终一致性。另一个代价是会拖慢应用程序的启动——除非你定期对状态做快照，或者能像上一节讨论的 CQRS 那样使用只读存储。

好了，CQRS 及相关模式就讲到这里。下面我们转向性能方面的另一个热门话题（此处不是双关）：缓存。

\subsection{缓存模式}

\label{idx:382}
恰当地使用缓存可以带来更好的性能、更低的延迟、更小的服务器负载（因而降低在云上运行的成本），并有助于缓解可扩展性方面的顾虑（所需的服务器更少）——有什么理由不喜欢呢？

\begin{myNotic}{注意}

如果你是来找关于 CPU 缓存的技巧的，可以在 \emph{第 13 章} \emph{性能}中找到。
\end{myNotic}

缓存是个很大的话题，我们在这里只涵盖其中的几个方面。

缓存的工作原理很简单：把读取最频繁的数据存放在访问速度快的非持久化存储中。缓存有许多不同的类型：

\begin{itemize}
\label{idx:383}
\item \textbf{客户端缓存}：用于专门为某个给定客户存储数据，通常放在客户端的机器或浏览器上。
\label{idx:384}
\item \textbf{Web 服务器缓存}：用于加速网页读取，例如通过 Varnish、NGINX 或 Apache 等能够缓存 Web 服务器响应的 HTTP 加速器。许多流行的网站和 Web 服务都会同时使用 NGINX 和 Varnish。
\label{idx:385}
\item \textbf{数据库缓存}：通过内置的、可调节的缓存机制来提升性能、可用性和可扩展性。
\label{idx:386}
\item \textbf{应用缓存}：用于加速你的应用，使它可以改为从缓存读取数据，而不是去访问它的数据库。
\item \textbf{CDN 也可以被当作缓存}：用于从靠近用户的位置提供内容，以降低延迟。
\end{itemize}

某些类型的缓存可以复制或部署为集群，以在规模化时提供性能。另一种选择是对它们进行分片：就像你会对数据库分片一样，你可以为数据的不同部分使用不同的缓存实例。

下面我们来逐一介绍更新缓存中数据的各种方法。毕竟，没人喜欢被提供陈旧的数据。

\subsubsection{更新缓存}

\label{idx:387}
有几种策略可以让缓存的数据保持新鲜。无论决定如何更新缓存项的是你自己还是另一家公司，了解这些策略都是值得的。本节我们将讨论它们的优缺点。

\paragraph{读穿透缓存}

\label{idx:388}
在这种策略中，数据按需加载：如果缓存命中，就立即返回。否则，在缓存未命中时，先从数据库获取数据，再更新缓存并返回数据。这种缓存适用于读取繁重的工作负载，但其缺点是在预取时会有一定的延迟。

\myGraphic{0.672}{images/B18456_04_18.png}{图 4.18：读穿透缓存}

\paragraph{写穿透缓存}

\label{idx:389}
如果你需要强一致性，那么同步更新数据库和缓存对你来说是一种可行的做法。这种策略可以保护你免于数据丢失：如果数据对用户可见，就意味着它已经写入数据库。写穿透缓存的缺点是，执行更新时的延迟比其他写入策略更大。

\myGraphic{0.686}{images/B18456_04_19.png}{图 4.19：写穿透缓存}

\paragraph{写回缓存}

\label{idx:390}
\label{idx:391}
另一种策略也称为\textbf{写回}，就是只给用户提供对缓存的访问。当用户执行更新时，缓存会把传入的更新排入队列，随后异步执行，从而更新数据库。这里显而易见的缺点是，如果出了什么差错，数据就永远无法写入。它实现起来也不像其他策略那么容易。不过它的好处是，用户所能感受到的是低延迟和高吞吐量。

\myGraphic{0.672}{images/B18456_04_20.png}{图 4.20：写回缓存}

\paragraph{旁路缓存}

\label{idx:392}
这种策略也称为\textbf{延迟加载}，其做法是按需填充缓存。它与读穿透的区别在于，由应用程序自己获取数据、更新缓存，并且在缓存失效时可以将其停用。

\myGraphic{0.531}{images/B18456_04_21.png}{图 4.21：旁路缓存}

这类缓存通常使用 Memcached、Valkey 或 Redis 来实现。它可以非常快速高效——缓存中只包含被请求过的数据。

\label{idx:393}
然而，如果经常请求不在缓存中的数据，前面提到的三次调用会明显增加延迟。为了在缓存重启时缓解这个问题，可以用持久化存储中选定的数据来为缓存预热（初始化）。

\paragraph{绕写缓存}

\label{idx:394}
在这种策略中，应用程序绕过缓存直接对数据库执行写操作，只有后续的读操作才会触发缓存。这种缓存适用于写入繁重的工作负载，且不会污染缓存，但其缺点是存在数据陈旧的风险，并且预取时会有一定的延迟。

\myGraphic{0.486}{images/B18456_04_22.png}{图 4.22：绕写缓存}

\paragraph{预刷新缓存}

\label{idx:395}
预刷新缓存会在热数据过期之前异步且周期性地对其进行更新，但其缺点是会额外增加缓存和数据库的负载。

\myGraphic{0.652}{images/B18456_04_23.png}{图 4.23：预刷新缓存}

\label{idx:396}
缓存中的条目可能会变得陈旧，因此最好为每个条目设置一个\textbf{生存时间 }（\textbf{TTL}）。如果要更新数据，可以采用写穿透的方式，把记录从缓存中移除并在数据库中更新它。使用仅基于时间的更新策略的多级缓存（例如 DNS 缓存）时要小心。这可能导致长时间使用陈旧数据，因为在 TTL 到期之前缓存可能都不会更新。

我们已经讨论了缓存的类型以及更新它们的策略，缓存就先讲到这里。下面我们转向构建可扩展架构的另一个方面：设计高效的数据存储。

\section{设计数据存储}

\label{idx:397}
下面我们来讨论应用程序的存储。首先，我们来决定应该采用 SQL、NoSQL 还是别的什么。

\label{idx:398}
大致来说，一个不错的经验法则是：根据数据库的规模来决定所用技术。对于小型数据库，比如说那些规模永远不会增长到 TB 级别的数据库，采用 SQL 是一种可行的做法。如果你的数据库非常小，或者想创建一个内存缓存，可以试试 SQLite。如果你打算进入个位数的 TB 范围，同样要保证规模永远不会超过这个范围，那么你最好的选择是采用 NoSQL。在某些情况下，仍然可以坚持使用 SQL 数据库，但由于硬件成本的原因，费用会迅速攀升，因为你的主节点需要一台性能强悍的服务器。即使这不是问题，你也应该衡量性能是否足以满足需求，并做好应对长时间维护窗口的准备。在某些情况下，使用 MySQL NDB Cluster 或 Citus（后者本质上是一个分片的 PostgreSQL）之类的技术来运行一组 SQL 机器集群，也可能适合你。然而，通常在这种情况下采用 NoSQL 会更便宜也更简单。如果你的数据库规模超过 10 TB，或者需要实时摄入数据，可以考虑改用数据仓库而不是 NoSQL。

\subsection{我应该使用哪种 NoSQL 技术？}

这个问题的答案取决于几个因素。这里列出了其中几个：

\begin{itemize}
\label{idx:399}
\item 如果你想存储时间序列（以很小且固定的间隔保存增量），那么最好的选择是使用 InfluxDB 或 VictoriaMetrics。
\item 如果你需要类似 SQL 的东西，但可以不要连接（join），换句话说，如果你打算把数据按列存储，可以试试 Apache Cassandra、ScyllaDB、AWS DynamoDB 或 Google 的 BigTable。
\label{idx:400}
\item 如果不是这种情况，那么你应该考虑你的数据是否是无模式的文档，例如 JSON 或某种应用程序日志。如果是，你可以采用 Elasticsearch，它非常适合这种灵活的数据，并提供 RESTful API。你也可以试试 RethinkDB 或 MongoDB，它们以\textbf{二进制 JSON（BSON）}格式存储数据，并支持 MapReduce。
\end{itemize}

好，但如果你不想存储文档呢？那么你可以选择对象存储，尤其是在数据量很大的情况下。通常在这种情况下选择云服务商是没问题的，这意味着使用 Amazon 的 S3、MinIO、Google Cloud Storage 或 Microsoft Azure Blob Storage 应该能帮上忙。如果你想采用本地方案，可以使用 OpenStack 的 Swift 或部署 Ceph。

如果文件存储也不是你想要的，那么也许你的场景只是简单的键值数据。使用这类存储有其好处，因为它很快。这正是许多分布式缓存基于它构建的原因。值得关注的技术包括 Riak、Redis、Valkey 和 Memcached（最后这个不适合持久化数据）。

除了前面提到的选项，你还可以考虑使用像 BerkeleyDB 这样的基于树的数据库。这类数据库本质上就是支持类似路径访问的特化键值存储。如果树对于你的场景限制太多，你可能会对 Neo4j 或 OrientDB 这类面向图的数据库感兴趣。

在理清数据存储之后，构建可扩展架构的另一个重要方面是规划如何将系统部署到生产环境。

\section{部署系统}

\label{idx:401}
尽管部署服务听起来很简单，但仔细一看，需要考虑的事情很多。本节将介绍如何高效地执行部署、在安装服务之后对其进行配置，以及检查它们在部署后是否保持健康，同时尽量减少停机时间。

\subsection{边车模式}

\label{idx:402}
\label{idx:403}
还记得本章前面提到的 Envoy 吗？它是一个非常有用的工具，可以助力高效的应用程序开发。你不必把日志、监控或网络等基础设施服务嵌入应用程序，而是可以把 Envoy 代理与你的应用一起部署，就像边车会\emph{被部署}在摩托车旁一样。两者结合起来所能做的事情，远比没有这个搭档（这种模式的另一个叫法）的应用要多得多。

使用边车可以加快开发速度，因为它带来的许多功能原本需要由你的每个微服务各自独立开发。由于边车与你的应用程序相互独立，因此可以使用你认为最适合这项工作的任何编程语言来开发它。边车以及它所提供的全部功能，可以由一个独立的开发者团队来维护，并与你的主服务独立地进行更新。

由于边车就驻留在它所增强的应用旁边，因此可以使用本地方式进行进程间通信。通常这已经足够快，并且比从另一台主机通信要快得多，但请记住，它可能会导致较高的系统资源消耗。

即使你部署的是第三方服务，在它旁边部署你选定的边车仍然能带来价值：你可以监控主机和服务的资源使用情况与状态，还可以在整个分布式系统中追踪请求。有时还可以根据服务的状况，通过编辑配置文件或 Web 界面来动态地重新配置该服务。

\label{idx:404}
在实际部署中，边车模式往往与其他架构模式搭配使用。例如，在 Docker 和 Kubernetes 环境中，边车模式通常与大使（ambassador）模式和适配器模式一起使用，用于收集、传输和分析容器的遥测数据以及进行通信。这些模式合在一起被称为三件套。\emph{图 4.24} 展示了这三者是如何协同工作的。

\myGraphic{0.920}{images/B18456_04_24.png}{图 4.24：大使、适配器与边车模式}

\label{idx:405}
\label{idx:406}
\label{idx:407}
\textbf{大使}是一个充当系统其他部分代理的容器；\textbf{适配器}是一个转换主容器输出的容器，比如为监控系统转换日志；\textbf{边车}则是一个增强或扩展主容器功能的额外容器。

这类组合部署模式往往是构建健壮、可投入生产的系统所必需的。

下面来看看使用 Envoy 实现边车（sidecar）模式的一个实践示例。

\subsubsection{使用 Envoy 部署带追踪与反向代理的服务}

\label{idx:408}
\label{idx:409}
我们将把 Envoy 用作部署的前端代理。在我们的配置中，Envoy 作为前端应用的代理，会把遥测数据发送给 OpenTelemetry Collector。首先创建 Envoy 的配置文件（本例中命名为 \texttt{envoy-front\_\allowbreak{}proxy.\allowbreak{}yaml}），并填入我们的代理地址：

\begin{codebox}[yaml]
static_resources:
  listeners:
  - address:
      socket_address:
        address: 0.0.0.0
        port_value: 8080
    traffic_direction: INBOUND
\end{codebox}

我们已经指定 Envoy 在端口 \texttt{8080} 上监听传入流量。稍后在配置中，我们会把它路由到我们的服务。现在，来指定我们希望使用一组服务实例来处理 HTTP 请求，并在其上加入一些追踪能力。首先，添加一个 HTTP 端点：

\begin{codebox}[yaml]
    filter_chains:
      - filters:
          - name: envoy.filters.network.http_connection_manager
            typed_config:
              "@type": type.googleapis.com/envoy.extensions.filters.network.http_connection_manager.v3.HttpConnectionManager
\end{codebox}

现在，来指定为请求分配 ID，并由分布式追踪系统 OpenTelemetry 进行追踪：

\begin{codebox}[yaml]
              generate_request_id: true
                tracing:
                  provider:
                    name: envoy.tracers.opentelemetry
                    typed_config:
                      "@type": type.googleapis.com/envoy.config.trace.v3.OpenTelemetryConfig
                      grpc_service:
                        envoy_grpc:
                          cluster_name: opentelemetry_collector
                        timeout: 0.250s
                      service_name: front_proxy
\end{codebox}

\label{idx:410}
我们将为请求创建 ID，并使用 OpenTelemetry 追踪器。该追踪器会把数据上报给运行在指定地址上的 OpenTelemetry Collector 实例。

\label{idx:411}
我们还需要指定，来自所有域的所有流量（参见 \texttt{match} 部分）都应路由到我们的服务集群：

\begin{codebox}[yaml]
              codec_type: auto
              stat_prefix: ingress_http
              route_config:
                name: example_route
                virtual_hosts:
                  - name: front_proxy
                    domains:
                      - "*"
                    routes:
                      - match:
                          prefix: "/"
                        route:
                          cluster: example_service
                        decorator:
                          operation: example_operation
\end{codebox}

稍后我们会定义 \texttt{example\_\allowbreak{}service} 集群。注意，每个到达该集群的请求都会被一个预定义的操作装饰器标记。我们还需要指定使用哪个路由器地址：

\begin{codebox}[python]
            http_filters:
            - name: envoy.filters.http.router
              typed_config:
                @type': type.googleapis.com/envoy.extensions.filters.http.router.v3.Router
            use_remote_address: true
\end{codebox}

\label{idx:412}
\label{idx:413}
现在我们知道了如何处理和追踪请求，剩下的就是定义我们所用到的集群。先从服务集群开始：

\begin{codebox}[yaml]
  clusters:
    - name: example_service
      connect_timeout: 0.250s
      type: strict_dns
      lb_policy: round_robin
      load_assignment:
        cluster_name: example_service
        endpoints:
          - lb_endpoints:
              - endpoint:
                  address:
                    socket_address:
                      address: example_service
                      port_value: 5678
    - name: opentelemetry_collector
      type: STRICT_DNS
      lb_policy: ROUND_ROBIN
      typed_extension_protocol_options:
        envoy.extensions.upstreams.http.v3.HttpProtocolOptions:
          "@type": type.googleapis.com/envoy.extensions.upstreams.http.v3.HttpProtocolOptions
          explicit_http_config:
            http2_protocol_options: {}
      load_assignment:
        cluster_name: opentelemetry_collector
        endpoints:
          - lb_endpoints:
              - endpoint:
                  address:
                    socket_address:
                      address: opentelemetry
                      port_value: 4317
\end{codebox}

\label{idx:414}
每个集群都可以包含我们服务的多个实例（端点）。在这里，如果我们决定添加更多端点，传入请求会使用轮询（round-robin）策略进行负载均衡。

再添加一个 \texttt{admin} 接口：

\begin{codebox}[yaml]
admin:
  access_log_path: /tmp/admin_access.log
  address:
    socket_address:
      address: 0.0.0.0
      port_value: 9901
\end{codebox}

现在把配置放进一个容器中，该容器会通过 Dockerfile 运行 Envoy，我们把它命名为 \texttt{Dockerfile-front\_\allowbreak{}proxy}：

\begin{codebox}[docker]
FROM envoyproxy/envoy:v1.32-latest
COPY --chmod=644 envoy-front_proxy.yaml /etc/envoy/envoy.yaml
\end{codebox}

\label{idx:415}
现在，来指定如何使用 Docker Compose 在多个容器中运行我们的代码。创建一个 Docker Compose（\texttt{compose.\allowbreak{}yaml}）文件，从定义前端代理服务开始：

\begin{codebox}[yaml]
services:
  front_proxy:
    build:
      context: .
      dockerfile: Dockerfile-front_proxy
    networks:
      - example_network
    ports:
      - "12345:12345"
      - "
9901:9901"
\end{codebox}

这里我们使用自己的 Dockerfile、一个简单的网络，并在宿主机上暴露容器中的两个端口：我们的服务和 \texttt{admin} 接口。现在添加代理将要指向的服务：

\begin{codebox}[yaml]
  example_service:
    image: hashicorp/http-echo
    networks:
      - example_network
    command: -text "It works!"
\end{codebox}

在本例中，该服务只会在一个简单的 Web 服务器中显示一个预定义的字符串。

\label{idx:416}
现在，我们在另一个容器中运行 OpenTelemetry Collector，并把它的端口暴露给外部：

\begin{codebox}[yaml]
  opentelemetry-collector:
    image: otel/opentelemetry-collector-contrib
    networks:
      - example_network
    ports:
      - "4317:4317"
      - "4318:4318"
      - "55679:55679"
    volumes:
      - ./collector-config.yaml:/etc/otelcol/config.yaml
\end{codebox}

最后一步是定义我们的网络：

\begin{codebox}[yaml]
networks:
  example_network: {}
\end{codebox}

\label{idx:417}
这样就完成了。现在你可以使用 \texttt{docker compose up --build} 运行该服务，并在浏览器中访问我们指定的端点。

使用边车代理有一个好处：即使你的服务挂掉了，边车通常仍然存活，可以在主服务宕机期间响应外部请求。当服务被重新部署时（例如因为一次更新）也是如此。说到这个，下面来学习如何最大限度地减少相关的停机时间。

\subsection{零停机部署}

\label{idx:418}
有两种常见方式可以最大限度地降低部署期间的停机风险：蓝绿部署和金丝雀发布。引入其中任何一种时，都可以使用 Envoy 边车。

\subsubsection{蓝绿部署模式}

\label{idx:419}
\textbf{蓝绿部署} 可以帮助你同时降低停机时间和与部署应用相关的回归风险（出现缺陷）。为此，你需要两个完全相同的生产环境，分别称为 \emph{蓝} 和 \emph{绿}。

\label{idx:420}
当绿色环境为顾客提供服务时，你可以在蓝色环境中执行更新。更新完成、服务经过测试、一切看起来都稳定之后，就可以切换流量，使其流向已更新的（蓝色）环境。

\label{idx:421}
如果切换后在蓝色环境中发现了任何问题，绿色环境仍然存在——你只需再切换回来即可。用户很可能根本注意不到任何变化，而且由于两个环境都在运行，切换过程中应该看不到任何停机。

\myGraphic{0.829}{images/B18456_04_25.png}{图 4.25：蓝绿部署}

\label{idx:422}
只需确保在切换过程中不会丢失任何数据（例如在新环境中进行的事务）。

\subsubsection{金丝雀发布模式}

\label{idx:423}
\label{idx:424}
要避免所有服务实例在更新后一起出故障，最简单的办法往往就是——不要一次更新全部实例。这正是蓝绿部署的渐进式变体（也称为金丝雀发布）背后的核心思想。

\myGraphic{0.812}{images/B18456_04_26.png}{图 4.26：金丝雀发布}

在 Envoy 中，你可以把下面的内容放进配置的 \texttt{routes} 部分：

\begin{codebox}[yaml]
- match:
    prefix: "/"
  route:
    weighted_clusters:
      clusters:
      - name: new_version
        weight: 5
      - name: old_version
        weight: 95
\end{codebox}

\label{idx:425}
\label{idx:426}
还要记得定义前面代码片段中提到的两个集群。下面是第一个，使用你服务的旧版本：

\begin{codebox}[yaml]
clusters:
  - name: old_version
    connect_timeout: 0.250s
    type: strict_dns
    lb_policy: round_robin
    load_assignment:
      cluster_name: old_version
      endpoints:
        - lb_endpoints:
            - endpoint:
                address:
                  socket_address:
                    address: old_version
                    port_value: 5678
\end{codebox}

第二个集群将运行新版本：

\begin{codebox}[yaml]
  - name: new_version
    connect_timeout: 0.250s
    type: strict_dns
    lb_policy: round_robin
    load_assignment:
      cluster_name: new_version
      endpoints:
        - lb_endpoints:
            - endpoint:
                address:
                  socket_address:
                    address: new_version
                    port_value: 5678
\end{codebox}

部署更新时，服务的新版本只会被一小部分用户（这里为 5\%）看到和使用。如果更新后的实例保持稳定，且没有任何检查与验证失败，你就可以分几步逐步更新更多主机，直到全部切换到新版本。你既可以手动更新配置文件，也可以使用 \texttt{admin} 端点来完成。大功告成！

\label{idx:427}
下面我们接着介绍本书要讲到的最后一个部署模式。

\subsection{外部配置存储模式}

\label{idx:428}
如果你部署的是一个简单应用，把配置和它一起部署也无妨。然而，当你希望进行更复杂的部署、且拥有大量应用实例时，仅仅为了重新配置就重新部署一版新应用，很快就会成为负担。与此同时，如果你想把服务当作牛群（cattle）而不是宠物（pet）来对待，手动修改配置是不可取的。引入外部配置存储是克服这些障碍的一种优雅方式。

本质上，你的应用可以从上述存储中获取配置，而不必仅依赖本地配置文件。这让你可以为多个实例提供公共设置、为其中部分实例调整参数，同时拥有一种简单而集中的方式来监控所有配置。如果你需要一个仲裁者来决定哪些节点是主节点、哪些作为备份节点，外部配置存储可以为实例提供这类信息。实现一套配置更新流程也很有用，这样实例在运行期间就能方便地重新配置。你可以使用现成的方案，如 Firebase Remote Config；利用基于 Java 的 Netflix Archaius；或者借助云存储和变更通知，自己编写一个配置存储。

\myGraphic{0.920}{images/B18456_04_27.png}{图 4.27：外部配置存储}

既然已经学习了一些有用的部署模式，下面我们来讨论高层设计中另一个重要主题：API。

\section{管理 API}

\label{idx:429}
设计良好的 API 对开发团队和产品的成功至关重要。我们可以把这个主题分为两类：系统级 API 和组件级 API。本节将讨论第一类层面上的 API 处理，下一章则会给出关于第二类的建议。

\subsection{API 管理工具}

\label{idx:430}
\label{idx:431}
除了管理对象，你还会希望管理整个 API。如果你想引入有关 API 使用的策略、控制对该 API 的访问、收集性能指标和其他分析数据，或者只是根据客户对你接口的使用情况进行收费，那么 \textbf{API 管理}（\textbf{APIM}）正是你要找的解决方案。

通常，一套 APIM 工具由以下组件组成：

\begin{itemize}
\item \textbf{API 网关}：API 所有用户的单一入口点。下一节会详细介绍。
\item \textbf{报告与分析}：用于监控 API 的性能和延迟、消耗的资源或发送的数据。这类工具可用于发现使用趋势，了解 API 的哪些部分及其背后的哪些组件是性能瓶颈，或者提供什么样的 SLA 才算合理，以及如何改进它们。
\item \textbf{开发者门户}：帮助开发者快速熟悉你的 API，并支持他们订阅你的 API。
\item \textbf{管理员门户}：用于管理策略和用户，并把 API 打包成可销售的产品。
\item \textbf{变现}：根据客户使用 API 的方式对其进行收费，并辅助相关的业务流程。
\end{itemize}

\label{idx:432}
APIM 工具既有云服务商提供的，也有独立第三方提供的，例如 NGINX 的 Controller 或 Tyk。

\label{idx:433}
为某个云平台设计 API 时，要了解该云服务商通常记录在文档中的最佳实践。例如，你可以在 \emph{延伸阅读} 一节找到 Google Cloud Platform 的常见设计模式。在它的场景中，很多实践都围绕使用 Protobuf 展开。

选择正确的 API 消费方式会大有裨益。向服务器发起请求最简单的方式是直接连接各个服务。这种方式易于搭建，对小型应用来说尚可，但日后可能引发性能问题。API 消费者很可能需要调用好几个不同的服务，从而导致高延迟。采用这种方式也无法实现真正的可扩展性。

一种更好的方法是使用 API 网关。这类网关通常是 APIM 解决方案不可或缺的一部分，但也可以单独使用。

\subsection{API 网关}

\label{idx:434}
API 网关是希望使用你 API 的客户端的入口点。它可以把传入请求路由到某个特定的服务实例或服务集群。这可以简化客户端代码，因为客户端不再需要了解所有后端节点或它们之间如何协作。客户端只需知道 API 网关的地址，其余都由网关处理。由于后端架构对客户端隐藏，因此可以轻松改造它，而无需触碰客户端代码。

\myGraphic{0.721}{images/B18456_04_28.png}{图 4.28：API 网关是 API 调用的入口点}

网关可以把系统 API 的多个部分聚合到单一入口点，并基于 URL、HTTP 头、Cookie 或消息内容等详细信息，使用 \textbf{七层路由} 把请求导向系统相应的部分。它在 \textbf{开放系统互连}（\textbf{OSI}）模型的应用层上运作，与之相对的是 \textbf{四层路由}，后者使用协议和端口等传输层细节，以及 IP 地址等网络层信息。

\myGraphic{0.540}{images/B18456_04_29.png}{图 4.29：OSI 模型的七层}

与本章描述的许多模式一样，始终要考虑为架构引入又一个模式而增加复杂度是否值得。想一想，如果可用性、容错性和性能对你很重要，添加它会对这些方面产生怎样的影响。毕竟，网关通常只是一个单节点，所以尽量不要让它成为瓶颈或单点故障。

我们在 \emph{第 2 章} \emph{架构风格} 中提到过 \emph{服务于前端的后端} 模式，它可以被视为 API 网关模式的一种变体。在 \emph{服务于前端的后端} 的场景中，每个前端都连接到各自的网关。

\section{小结}

在本章中，我们学到了不少东西。现在你知道了何时该应用哪种服务模型，以及如何避免设计分布式系统时的常见陷阱。你了解了 CAP 和 PACELC 定理，以及它们对分布式架构产生的实际影响。现在你能够在这类系统中成功运行事务、降低停机时间、预防问题，并从容地从错误中恢复。应对异常高的负载不再是黑魔法。把系统的各个部分（甚至是遗留部分）与你新设计的部分集成起来，你也能做到了。此外，你还掌握了一些提高系统性能和可扩展性、以及设计数据存储的技巧。部署系统和对系统进行负载均衡也不再神秘，现在你可以高效地完成它们。最后同样重要的是，发现服务以及设计和管理它们的 API，这些你现在都已学会如何去做。不错！

在下一章中，我们将学习如何利用具体的 C++ 特性，以更轻松、更高效的方式踏上通往卓越架构之路。

\section{问题}

\begin{enumerate}
\item 什么是事件溯源？
\item CAP 和 PACELC 定理常见的实际影响是什么？
\item Netflix 的 Chaos Monkey 可以用来做什么？
\item 缓存可以应用在哪些地方？
\item 当整个数据中心宕机时，你如何防止应用也随之下线？
\item 为什么要使用 API 网关？
\item Envoy 如何帮助你实现各种架构目标？
\end{enumerate}

\section{延伸阅读}

\begin{itemize}
\item Microsoft Azure，\emph{Cloud Design Patterns}：\href{https://docs.microsoft.com/en-us/azure/architecture/patterns/ }{https://docs.microsoft.com/en-us/azure/architecture/patterns/}
\item Google API Improvement Proposals，\emph{Design Patterns}：\href{https://google.aip.dev/general\#design-patterns }{https://google.aip.dev/general\#design-patterns}
\item Microsoft，\emph{Microsoft REST API Guidelines}，GitHub：\href{https://github.com/microsoft/api-guidelines/blob/vNext/Guidelines.md}{https://github.com/microsoft/api-guidelines/blob/vNext/Guidelines.md}
\item Envoy，\emph{Getting Started}：\href{https://www.envoyproxy.io/docs/envoy/latest/start/start }{https://www.envoyproxy.io/docs/envoy/latest/start/start}
\item Jay Runkel，\emph{Active-Active Application Architectures with MongoDB}（博客），MongoDB：\href{https://www.mongodb.com/blog/post/active-active-application-architectures-with-mongodb}{https://www.mongodb.com/blog/post/active-active-application-architectures-with-mongodb}
\item DB-Engines，\emph{DB-Engines Ranking}：\href{https://db-engines.com/en/ranking }{https://db-engines.com/en/ranking}
\item Jasper，\emph{How Millisecond Delays May Kill Database Performance}（博客），Packet-Foo：\href{https://blog.packet-foo.com/2014/09/how-millisecond-delays-may-kill-database-performance/ }{https://blog.packet-foo.com/2014/09/how-millisecond-delays-may-kill-database-performance/}
\item Istio，\emph{Architecture}：\href{https://istio.io/latest/docs/ops/deployment/architecture/ }{https://istio.io/latest/docs/ops/deployment/architecture/}
\end{itemize}

\begin{center}
\begin{tabularx}{\textwidth}{XX}
\toprule
\textbf{立即解锁本书的专属福利} 扫描此二维码，或访问 \href{http://packtpub.com/unlock}{packtpub.com/unlock}，然后按书名搜索本书。 & \begin{center}\includegraphics[width=0.250\linewidth]{images/Unlock_Now.png}\end{center} \\
\midrule
\emph{注意：开始之前请准备好你的购买发票。} &  \\
\bottomrule
\end{tabularx}
\end{center}
